\documentclass[11pt]{article}

\usepackage[margin=1in]{geometry}
\usepackage{fontspec}
\setmainfont{TeX Gyre Termes}
\setmonofont{TeX Gyre Cursor}[Scale=0.88]
\usepackage{amsmath,amssymb}
\usepackage{booktabs}
\usepackage{array}
\usepackage{xcolor}
\usepackage{listings}
\usepackage{enumitem}
\usepackage{hyperref}
\usepackage{caption}
\usepackage{fancyhdr}
\usepackage{titlesec}
\usepackage{longtable}
\usepackage{multirow}

\hypersetup{
  colorlinks=true,
  linkcolor=black,
  citecolor=black,
  urlcolor=blue,
  pdftitle={Cryptanalysis of KOD8\_ENC (KOD8 v2.0)},
  pdfauthor={Fatih}
}

\definecolor{codebg}{RGB}{245,245,245}
\definecolor{codekw}{RGB}{0,90,150}
\definecolor{codestr}{RGB}{140,20,20}
\definecolor{codecomm}{RGB}{110,110,110}

\lstset{
  backgroundcolor=\color{codebg},
  basicstyle=\ttfamily\small,
  keywordstyle=\color{codekw}\bfseries,
  commentstyle=\color{codecomm}\itshape,
  stringstyle=\color{codestr},
  showstringspaces=false,
  breaklines=true,
  breakatwhitespace=false,
  frame=single,
  framesep=4pt,
  language=Python,
  numbers=left,
  numberstyle=\tiny\color{gray},
  numbersep=6pt,
  xleftmargin=1.2em,
  tabsize=4
}

\lstdefinestyle{appendixcode}{
  backgroundcolor=\color{codebg},
  basicstyle=\ttfamily\footnotesize,
  keywordstyle=\color{codekw}\bfseries,
  commentstyle=\color{codecomm}\itshape,
  stringstyle=\color{codestr},
  showstringspaces=false,
  breaklines=true,
  breakatwhitespace=false,
  frame=single,
  framesep=4pt,
  language=Python,
  numbers=left,
  numberstyle=\tiny\color{gray},
  numbersep=6pt,
  xleftmargin=1.4em,
  tabsize=4
}

\pagestyle{fancy}
\fancyhf{}
\fancyhead[L]{\small Cryptanalysis of KOD8 v2.0}
\fancyhead[R]{\small \thepage}
\renewcommand{\headrulewidth}{0.3pt}

\titleformat{\section}{\normalfont\Large\bfseries}{\thesection}{1em}{}
\titleformat{\subsection}{\normalfont\large\bfseries}{\thesubsection}{1em}{}
\titleformat{\subsubsection}{\normalfont\normalsize\bfseries}{\thesubsubsection}{1em}{}

\title{\textbf{Cryptanalysis of KOD8 Encryption Method:\\
A Claims-Traceable Audit of a Custom Eight-Stage Symmetric Cipher}}
\author{A. Fatih ULUSOY}
\date{August 2026}

\begin{document}

\maketitle

\begin{abstract}
\noindent
KOD8 Encryption Method is a custom symmetric cipher, described in
its own project documentation as a ``multi-layered obfuscation and
encryption engine,'' that routes input through one of six
input-type-specific ``CipherLists,'' each composing exactly eight
primitive reversible transformations (encodings, XOR variants, classical
substitutions, and transpositions) driven by a single eight-digit numeric
key. This report presents a structural and empirical cryptanalysis of
the reference implementation, structured explicitly as a
claims-versus-findings audit against the project's own published
feature list, combining static source analysis, worked
per-CipherList traces, avalanche-criterion measurements, ciphertext
randomness testing, a full classical frequency-analysis attack, and a
benchmarked known-plaintext brute-force key recovery. We show that the
cipher's advertised eight-layer, ``high-diffusion'' design provides far
less effective key-dependent entropy than its layer count suggests: six
of seventeen primitives never reference the key at all, one collapses
an eight-digit (nominally $10^{8}$) key into at most 325 distinct
internal states, and the block-cipher-inspired CBC-style layer reuses
the literal key string as a fixed, non-random initialization vector
across every message directly contradicting the specific
per-CipherList claims made for CL3 and CL5 in the project's own
documentation. We measure the avalanche effect directly and find it far
below the ideal 50\% criterion for every CipherList (as low as 0.1\% for a single-digit
key change under CL1), demonstrate that the cipher's own ciphertext has
\emph{lower} Shannon entropy (3.31 bits/byte) than its plaintext input
(4.19 bits/byte), and the central empirical result of this report ---
show that a purely ciphertext-only chi-square frequency attack recovers
\textbf{all eight digits} of the master key from the Vigen\`ere-style
substitution layer alone, with no brute force and no known plaintext
required. We further confirm the cipher is fully deterministic --- the
project's own README names this property ``Deterministic Security,''
which Section~\ref{sec:determinism} shows is a contradiction in terms ---
is unauthenticated (several primitives silently absorb decoding
errors rather than signal tampering), and that the reference application
persists the selected CipherList identifier as a cleartext header,
eliminating any classification uncertainty an attacker would otherwise
face. A conservative, no-shortcuts known-plaintext brute-force attack
over the full nominal keyspace recovers the key for the plain-text
CipherList in under 2.4 hours on a single CPU core using an unoptimized
reference Python implementation. Despite these findings, we note that
the project's own license text already scopes KOD8\_ENC for
``educational and obfuscation purposes'' --- precisely the classification
this report independently arrives at --- and a full claims traceability
matrix (Table~\ref{tab:traceability}) is provided so that each
individual marketing claim can be checked against its corresponding
evidence. We conclude that KOD8 is correctly classified as a composite
classical (product) cipher and a functional obfuscation layer, not as a
cryptographically secure primitive, and we outline concrete
remediations.
\end{abstract}

\tableofcontents
\newpage

% =====================================================================
\section{Executive Summary}
% =====================================================================

This report answers a single practical question: \emph{if an adversary
had to break KOD8\_ENC, how would they do it, and how long would it
take?} We answer it five ways: by reading the project's own stated
claims and holding them to account, by reading the source line by
line, by measuring the cipher's statistical behavior directly, by
mounting a textbook 19th-century frequency attack against one of its
layers, and by running an actual brute-force program against the
reference implementation.

All five approaches agree. The headline results are:

\begin{itemize}[leftmargin=1.6em]
\item \textbf{The project's own README already scopes it correctly}:
its license section states plainly that KOD8\_ENC is ``provided for
educational and obfuscation purposes.'' Every finding below is
consistent with that self-assessment; where this report disagrees with
the README, it is with specific per-feature marketing claims
(``high-diffusion,'' ``Deterministic Security,'' individual
CipherList taglines), not with the project's own bottom-line framing.
A full claim-by-claim traceability matrix appears in
Table~\ref{tab:traceability}.
\item \textbf{Six of the cipher's seventeen primitive operations never
read the key at all.} Once the algorithm is known which, per
Kerckhoffs's principle, must always be assumed these steps are free
to invert.
\item \textbf{The substitution-box layer collapses an eight-digit key
into 325 possible internal states}, due to a seeding formula that sums
weighted digits rather than hashing the full key.
\item \textbf{The CBC-style block layer reuses the literal key string
as its initialization vector}, on every message, forever directly
contradicting the README's claim that this layer exists to ``break
long repeating patterns.''
\item \textbf{The avalanche effect --- the standard measure of how well a
cipher scrambles small input changes --- is far below the accepted 50\%
ideal for every CipherList we tested}, and is effectively absent
(0.1\%) for a single-digit key change under the plain-text CipherList ---
directly at odds with the README's ``high-diffusion data protection''
claim.
\item \textbf{The cipher's own ciphertext is measurably less random
than its plaintext} (3.31 vs.\ 4.19 Shannon bits/byte), and fails a
basic chi-square uniformity test by four orders of magnitude.
\item \textbf{A purely ciphertext-only classical frequency attack
recovers the entire eight-digit master key} from the Vigen\`ere-style
layer, with no brute force and no known plaintext directly
contradicting the README's claim that this same layer exists to
``eliminate decimal patterns'' when reused as \texttt{op\_base36} in CL5.
\item \textbf{A conservative brute-force attack recovers the key in
under 2.4 hours} on a single core with zero optimization, and in
minutes to seconds once parallelized or once the structural shortcuts
above are exploited.
\end{itemize}

None of this makes KOD8\_ENC a bad \emph{learning exercise} --- it is a
genuinely interesting worked example of a hand-built product cipher,
and Section~\ref{sec:discussion} argues that, evaluated strictly as an
obfuscation layer against casual or automated pattern-matching
inspection, several of its ``zero-entropy'' steps do real work. But it
should not be relied upon to keep anything confidential against a
motivated reader of its own source code --- a conclusion the project's
own documentation, read carefully, already reaches.

% =====================================================================
\section{Introduction}
% =====================================================================

KOD8 (``Kod'' Turkish for ``code'') is a hand-designed symmetric
cipher implemented in \texttt{kod8\_engine.py}, with a companion
Tkinter front end, \texttt{kod8\_ui.py}, that visualizes the encryption
and decryption process step by step. The engine's own documentation
frames it as an ``anti-LLM and decoding agent based cipher model'' ---
that is, a scheme designed primarily to resist casual or automated
pattern-based decoding rather than to meet a formal security definition
such as IND-CPA or IND-CCA2.

This framing matters for how the design should be evaluated. KOD8 is
not positioned as a competitor to AES-GCM or ChaCha20-Poly1305; it is a
personal cryptographic-engineering exercise, and a legitimate,
instructive one: a working example of a \emph{product cipher} in the
classical sense --- a chain of substitution and transposition steps, in
the tradition of ciphers that predate the computational security
paradigm entirely.

The purpose of this report is to evaluate that claim rigorously, using
four complementary methods:

\begin{enumerate}[leftmargin=1.6em]
  \item \textbf{Static structural analysis} --- how much of the nominal
        $10^{8}$-key-space entropy actually reaches each of the eight
        steps in a given CipherList (Section~\ref{sec:entropy}),
        traced against every primitive's source.
  \item \textbf{Dynamic statistical analysis} --- does the cipher exhibit
        the avalanche effect and output randomness expected of a sound
        cipher, measured directly (Sections~\ref{sec:avalanche}
        and~\ref{sec:randomness})?
  \item \textbf{Classical cryptanalysis} --- can a 19th-century
        frequency-analysis technique, applied without any brute force
        or known plaintext, recover key material
        (Section~\ref{sec:freq})?
  \item \textbf{Empirical brute force} --- how long does exhaustive key
        recovery actually take against the reference implementation,
        measured rather than estimated (Section~\ref{sec:brute})?
\end{enumerate}

Every quantitative claim in this report was produced by executing the
actual reference engine, not by inspection alone; source excerpts are
included throughout, and full listings of both reference files appear
in the appendices.

% =====================================================================
\section{The Project's Own Description: KOD8\_ENC}
\label{sec:readme}
% =====================================================================

Before presenting our own analysis, we record the project's stated
purpose and design claims in its own words, as published in its
\texttt{README.md}. Doing so serves two purposes: it gives the reader
the intended design context this report is evaluating against, and it
lets us structure the remainder of the report as a direct
\emph{claims-versus-findings} audit rather than an unanchored critique.

\subsection{Stated Purpose and Feature Set}

The project identifies itself as \textbf{KOD8\_ENC} and describes
itself as follows:

\begin{quote}
\itshape
``KOD 8 is a multi-layered obfuscation and encryption engine designed
for high-diffusion data protection. Unlike standard single-algorithm
ciphers, KOD 8 utilizes a Strategic Chain Architecture, applying
exactly eight distinct, reversible transformations based on the
detected or specified input type.''
\end{quote}

The README lists six headline features, each of which we treat as an
individually verifiable claim in Table~\ref{tab:traceability}:
\emph{Strategic Auto-Detection}, \emph{8-Round Chain Logic},
\emph{Turkish Alphabet Support}, \emph{Zero-Dependency UI},
\emph{Deterministic Security}, and a documented set of \emph{Chain
Safety Rules (R1–R6)} intended to guarantee lossless reversibility.
Notably, the feature is named \emph{``Deterministic Security''} in the
project's own words --- Section~\ref{sec:determinism} shows that
determinism and semantic security are, in the standard cryptographic
sense, in direct tension with one another; the project's own naming
inadvertently identifies the property we go on to show is a
weakness, not a strength.

\subsection{The Conceptual Base Function}

The README offers a compact conceptual formula for the ``8-Round Chain
Logic'':

\begin{quote}
\itshape
$F(x) = 1(3(9(7(4(2(6(8 \cdot x)))))))$
\end{quote}

Read from the innermost term outward, the sequence of numeric labels
is $8, 6, 2, 4, 7, 9, 3, 1$ --- which is exactly the digit sequence of
the reference master key, \texttt{KOD8\_KEY = "86247931"}, read left
to right. This confirms that the formula is a conceptual mnemonic for
``eight rounds, ordered by the eight digits of the key'' rather than a
literal description of the arithmetic each primitive performs (the
real primitives are byte-level substitutions, transpositions, and XOR
operations, not scalar multiplication by a labeled constant, as
Section~\ref{sec:entropy} details). We note this explicitly because it
is a clean, independent confirmation derived from the project's own
marketing formula rather than from the source code of a fact central
to this report: \textbf{every stage of the chain is nominally
associated with one digit of the key, but, as Table~\ref{tab:entropy}
shows, most stages never actually consult the digit conceptually
assigned to them.}

\subsection{Stated Per-CipherList Focus}

The README assigns each CipherList a targeted data type and a
one-line ``primary focus,'' reproduced in Table~\ref{tab:readme-cl},
which we use throughout the remainder of this report as the
\emph{intended} security property of each chain, to be checked against
what the chain actually delivers.

\begin{table}[h]
\centering
\small
\caption{Per-CipherList targeted data and stated primary focus, as published in \texttt{README.md}.}
\label{tab:readme-cl}
\begin{tabular}{@{}lll@{}}
\toprule
\textbf{ID} & \textbf{Targeted data} & \textbf{Stated primary focus} \\
\midrule
CL1 & Human-readable prose & Turkish-aware substitution \& linguistic diffusion \\
CL2 & PNG, JPG, WebP & Binary-safe S-Box layers \& block shuffling \\
CL3 & MP4, MKV, AVI & Stream-XOR heavy to break long repeating patterns \\
CL4 & PDF, DOCX, TXT & Structure-breaking columnar transpositions \\
CL5 & JSON, CSV, coordinates & Base-36 transformation to eliminate decimal patterns \\
CL6 & Mixed/unknown & Maximum obfuscation using all primitive families \\
\bottomrule
\end{tabular}
\end{table}

\subsection{Chain Safety Rules and License Scope}

The README documents six ``Chain Safety Rules (R1–R6)'' whose
stated purpose is exclusively \emph{reversibility} --- ensuring
encryption and decryption round-trip without data corruption --- not
confidentiality. Two are called out explicitly: \emph{R4}, which
requires that \texttt{unicode\_shift} never precede Turkish-alphabet
operations, and \emph{R6}, which requires all text input be normalized
to uppercase before the first step. Our own analysis of every
CipherList's step ordering (Table~\ref{tab:chains},
Section~\ref{sec:traces}) is consistent with both rules being followed
correctly throughout the reference implementation: as an
\emph{engineering-correctness} property the chain does round-trip
losslessly --- KOD8\_ENC delivers exactly what it claims. This is a
materially different claim from a security property, and our findings
below should not be read as disputing it.

Finally, and importantly, the README's own license section states the
project's intended scope plainly:

\begin{quote}
\itshape
``This project is provided for educational and obfuscation purposes.''
\end{quote}

This is the single most important sentence in the source repository
for the purpose of this report. It is also the one piece of evidence
in this entire analysis that we did not have to derive ourselves: the
project already classifies itself in exactly the category
Section~\ref{sec:discussion} independently arrives at through
structural and empirical analysis. The remainder of this report can
therefore be read less as a rebuttal of KOD8\_ENC's claims and more as
a rigorous confirmation of the one claim that matters most, together
with a precise accounting of where the more specific marketing claims
(``high-diffusion,'' ``Deterministic Security,'' per-CipherList
``primary focus'' descriptions) overstate what the implementation
actually achieves against a cryptanalytic adversary as opposed to
against casual or automated pattern-based inspection, which is a
target the license text itself signals the project was never trying to
exceed.

\subsection{Claims Traceability Matrix}

Table~\ref{tab:traceability} maps every substantive claim in the
README against the section of this report that verifies it, and a
verdict: \textbf{Confirmed} (the implementation delivers what is
claimed), \textbf{Partially Confirmed} (true in a narrower or weaker
sense than the claim implies), or \textbf{Refuted} (the implementation
does not deliver the claimed property against the threat model in
Section~\ref{sec:threat}).

{\small
\begin{longtable}{@{}p{3.0cm}p{4.8cm}p{2.0cm}p{2.7cm}@{}}
\caption{Claims traceability matrix: every substantive claim in \texttt{README.md}, the report section that verifies it, and the verdict.}
\label{tab:traceability} \\
\toprule
\textbf{Claim (as stated)} & \textbf{Verification} & \textbf{Verdict} & \textbf{Evidence} \\
\midrule
\endfirsthead
\multicolumn{4}{l}{\textit{Table \thetable\ continued from previous page}} \\
\toprule
\textbf{Claim (as stated)} & \textbf{Verification} & \textbf{Verdict} & \textbf{Evidence} \\
\midrule
\endhead
\bottomrule
\endfoot
``8-Round Chain Logic'' / exactly 8 operations & Every CipherList does run exactly 8 steps & Confirmed & \S\ref{sec:traces}, App.~\ref{app:engine} \\
Each round meaningfully uses the key & Only 2–4 of 8 steps per chain reference the key at all & Refuted & \S\ref{sec:entropy}, Table~\ref{tab:entropy} \\
``High-diffusion data protection'' & Avalanche effect far below 50\% ideal for every chain; 0.1\% for CL1 key changes & Refuted & \S\ref{sec:avalanche}, Table~\ref{tab:avalanche} \\
``Deterministic Security'' & Determinism $\Rightarrow$ no semantic security; identical plaintexts always yield identical ciphertexts & Refuted & \S\ref{sec:determinism} \\
Strategic Auto-Detection (6 lists) & Chain identity is stored in a cleartext file header, eliminating ambiguity & Partially Confirmed & \S\ref{sec:cid} \\
Turkish Alphabet Support & Feature works as implemented; is also the primary attack surface & Partially Confirmed & \S\ref{sec:freq} \\
CL3: ``break long repeating patterns'' & Fixed, key-derived (non-random) CBC-style IV reused on every message & Refuted & \S\ref{sec:iv} \\
CL5: ``eliminate decimal patterns'' & Base-36/Vigen\`ere layer fully broken by classical frequency analysis & Refuted & \S\ref{sec:freq} \\
CL4: ``structure-breaking columnar transpositions'' & Columnar step has only 5 possible states regardless of key & Partially Confirmed & \S\ref{sec:entropy}(b) \\
Customizable master key & Key is a hardcoded source-file constant; no runtime key input exists & Partially Confirmed & \S\ref{sec:threat} (T1) \\
Chain Safety Rules R1–R6 & Verified consistent with observed chain orderings & Confirmed & \S\ref{sec:readme}, \S\ref{sec:traces} \\
License: ``educational and obfuscation purposes'' & Matches this report's independently derived classification & Confirmed & \S\ref{sec:discussion} \\
\end{longtable}
}

% =====================================================================
\section{Background and Related Concepts}
% =====================================================================

This section briefly reviews the classical cryptographic concepts this
report relies on, for a reader encountering them for the first time
alongside KOD8 itself.

\subsection{Kerckhoffs's Principle}

A cipher's security should depend only on the secrecy of its key, never
on the secrecy of its algorithm. This nineteenth-century design
principle (Kerckhoffs, 1883) is the standard baseline assumption in
modern cryptanalysis: an evaluator assumes the attacker already knows
exactly how the cipher works, and asks what remains protected once that
assumption is granted. Every ``zero-entropy'' finding in
Section~\ref{sec:entropy} follows directly from applying this
principle to KOD8's primitives.

\subsection{Product Ciphers, Confusion, and Diffusion}

Shannon's 1949 treatment of secrecy systems introduced two properties a
strong cipher should exhibit: \emph{confusion} (the relationship
between key and ciphertext should be as complex as possible --- typically
achieved via substitution) and \emph{diffusion} (each plaintext or key
bit should influence many ciphertext bits --- typically achieved via
transposition or mixing). A \emph{product cipher} chains multiple
simple substitution and transposition operations together specifically
to compound these properties. KOD8's eight-step chains are, structurally,
product ciphers in exactly this sense; Section~\ref{sec:avalanche}
measures how well they actually achieve diffusion in practice.

\subsection{The Avalanche Criterion}
\label{sec:avalanche-criterion-bg}

A widely used practical test of diffusion is the \emph{avalanche
effect}: flipping a single input bit (of plaintext or key) should, on
average, flip approximately half of the output bits. This 50\% figure
is not an arbitrary target --- it is the expected behavior of an ideal
random function, and ciphers that fall well short of it leave
statistical structure an attacker can exploit. Section~\ref{sec:avalanche}
measures this directly against the reference KOD8 engine.

\subsection{Classical Frequency Cryptanalysis}

Long before computational cryptanalysis existed, repeating-key
substitution ciphers (of which the Vigen\`ere cipher is the canonical
example) were broken by exploiting the non-uniform frequency
distribution of natural-language letters. Once an attacker knows (or
guesses) the key period, the ciphertext can be split into one
sub-stream per key position; each sub-stream is then just a
single-shift substitution cipher, breakable by comparing its observed
letter-frequency distribution against a known reference distribution
(a \emph{chi-square goodness-of-fit} test is a standard way to score
candidate shifts). Section~\ref{sec:freq} applies this exact technique,
essentially unchanged since the nineteenth century, to KOD8's
Vigen\`ere-style layer.

\subsection{Linear-Congruential Generators as Keystreams}

A linear congruential generator (LCG) produces a pseudorandom sequence
via $s_{i+1} = (a \cdot s_i + c) \bmod m$ for fixed constants $a, c, m$.
LCGs are fast and simple but are well known to be cryptographically
weak: their internal state is small, their output is statistically
distinguishable from random, and critically for a keystream
application --- recovering a handful of consecutive outputs is often
sufficient to reconstruct the seed and predict all future (and past)
outputs. KOD8's keystream and S-box derivation both use an LCG with the
classic constants $a=1103515245$, $c=12345$ (the same constants used by
the glibc \texttt{rand()} implementation), which are public knowledge
independent of this report.

% =====================================================================
\section{System Overview}
% =====================================================================

KOD8 defines a single global key,

\begin{lstlisting}
KOD8_KEY = "86247931"
\end{lstlisting}

\sloppy
\noindent
an eight-character decimal string, and a fixed 29-letter Turkish
alphabet used by two of its substitution primitives:
\texttt{TURKISH\_ALPHA = "ABCÇDEFGĞHİIJKLMNOÖPRSŞTUÜVYZ"}. Seventeen
primitive operations are defined in total --- two encodings
(\texttt{op\_hex}, \texttt{op\_base64}), eight XOR/substitution
variants (\texttt{op\_xor\_key}, \texttt{op\_rolling\_xor},
\texttt{op\_keystream\_xor}, \texttt{op\_vigenere\_tr},
\texttt{op\_atbash\_tr}, \texttt{op\_unicode\_shift}, \texttt{op\_sbox},
\texttt{op\_base36}), and seven transposition/block operations
(\texttt{op\_rail\_fence}, \texttt{op\_block\_rotate},
\texttt{op\_split\_reverse}, \texttt{op\_full\_reverse},
\texttt{op\_columnar}, \texttt{op\_block\_shuffle},
\texttt{op\_block\_xor\_cbc}) --- each with signature
\texttt{op\_xxx(text, encrypt, **kwargs)} and each claimed to be
self-contained and reversible.
\fussy

\subsection{Formal Notation}

We write a CipherList as an ordered composition of eight primitive
functions $f_1, \dots, f_8$, each parameterized by the global key $K$
and, for some primitives, by fixed per-list literals $\ell$ (such as a
rail count or rotation width):

\[
  E_{\text{CL}}(P) = f_8(K,\ell_8) \circ f_7(K,\ell_7) \circ \cdots \circ f_1(K,\ell_1)\,(P)
\]

Decryption applies the inverse of each $f_i$ in reverse order:

\[
  D_{\text{CL}}(C) = f_1^{-1}(K,\ell_1) \circ f_2^{-1}(K,\ell_2) \circ \cdots \circ f_8^{-1}(K,\ell_8)\,(C)
\]

The engine additionally defines an auto-detector
$\delta: \text{input} \to \{\text{CL1},\dots,\text{CL6}\}$ that selects
which composition to apply, and the reference \texttt{Kod8.encrypt}
method returns the pair $(C, \text{cid})$ where $\text{cid} = \delta(P)$.
As shown in Section~\ref{sec:cid}, the reference UI persists
$\text{cid}$ in cleartext, so for the remainder of this report we treat
$\text{cid}$ as known to the attacker and focus entirely on recovering
$K$ for a single, identified chain.

\subsection{The Six CipherLists}

Six CipherLists (CL1–CL6) each select exactly eight of the primitives
above, in a fixed order, keyed to an input type inferred by a heuristic
auto-detector (\texttt{auto\_detect}): plain text, image, video,
document, numeric/data, and an ``experimental'' catch-all used when no
other detector matches. Table~\ref{tab:chains} lists the encrypt-order
chain for each.

\begin{table}[h]
\centering
\small
\caption{The eight-step encrypt-order chain defined for each CipherList.}
\label{tab:chains}
\begin{tabular}{@{}lp{10.5cm}@{}}
\toprule
\textbf{ID} & \textbf{Encrypt chain (left to right)} \\
\midrule
CL1 & hex $\to$ xor\_key $\to$ vigenere\_tr $\to$ rail\_fence(rails=3) $\to$ atbash\_tr $\to$ block\_rotate(n=3) $\to$ unicode\_shift $\to$ full\_reverse \\
CL2 & base64 $\to$ sbox $\to$ rolling\_xor $\to$ block\_shuffle $\to$ xor\_key $\to$ split\_reverse $\to$ hex $\to$ unicode\_shift \\
CL3 & base64 $\to$ keystream\_xor $\to$ block\_xor\_cbc $\to$ block\_shuffle $\to$ rolling\_xor $\to$ rail\_fence(rails=4) $\to$ split\_reverse $\to$ full\_reverse \\
CL4 & hex $\to$ columnar $\to$ vigenere\_tr $\to$ rail\_fence(rails=5) $\to$ block\_shuffle $\to$ xor\_key $\to$ atbash\_tr $\to$ split\_reverse \\
CL5 & base36 $\to$ hex $\to$ vigenere\_tr $\to$ block\_rotate(n=3) $\to$ sbox $\to$ keystream\_xor $\to$ block\_shuffle $\to$ full\_reverse \\
CL6 & hex $\to$ sbox $\to$ block\_xor\_cbc $\to$ columnar $\to$ block\_shuffle $\to$ keystream\_xor $\to$ atbash\_tr $\to$ full\_reverse \\
\bottomrule
\end{tabular}
\end{table}

\noindent
\emph{(Chain orders extracted programmatically from the live
\texttt{CIPHER\_LISTS} dictionary in \texttt{kod8\_engine.py} by
iterating each CipherList's step list and its bound keyword arguments,
rather than transcribed by hand, to avoid transcription error. Full
source appears in Appendix~\ref{app:engine}.)}

The Tkinter UI (\texttt{kod8\_ui.py}) wraps this engine, and critically
for the threat model developed in Section~\ref{sec:threat} persists
encrypted output to disk in a fixed text format:

\begin{lstlisting}[language=Python]
f.write(f"# KOD8  cipher_id: {cid}\n{enc}")
\end{lstlisting}

% =====================================================================
\section{Threat Model}
\label{sec:threat}
% =====================================================================

We evaluate KOD8 under three progressively weaker attacker assumptions,
each grounded in how the reference implementation actually behaves.

\begin{description}[leftmargin=1.4em, style=nextline]
  \item[\textbf{T1 Source-available (white-box).}]
    \texttt{KOD8\_KEY} is a hardcoded module-level constant, not a
    user-supplied secret, a password-derived key, or a per-deployment
    value. Any party with access to \texttt{kod8\_engine.py} --- including
    anyone who has ever received this repository, cloned it from a
    version-control host, or extracted it from a distributed build ---
    possesses the key outright. Under Kerckhoffs's principle this would
    normally be an acceptable assumption \emph{provided the key itself
    remained secret}; here the key \emph{is} the source file. This is the
    single most severe practical weakness in the system and requires no
    cryptanalysis at all to exploit.
  \item[\textbf{T2 Ciphertext-and-header-only, key rotated.}]
    A more charitable model assumes a deployer has changed
    \texttt{KOD8\_KEY} to a private eight-digit value not present in any
    public source tree, and that the attacker only observes
    \texttt{.kod8} files or wire traffic. Because the CipherList
    identifier is stored as a cleartext header (Section~\ref{sec:cid}),
    the attacker's search problem is immediately reduced to: recover one
    eight-digit key against one \emph{known} 8-step chain --- not, as the
    ``six independent CipherLists'' design might suggest, a joint
    key-and-chain identification problem.
  \item[\textbf{T3 Known-plaintext / ciphertext-only.}]
    The strongest assumptions we exercise empirically in this report are
    (a) a single known plaintext/ciphertext pair (Section~\ref{sec:brute}),
    which is a routine assumption in classical cryptanalysis and
    realistic here given KOD8's own self-test harness, UI placeholder
    text, and predictable file headers; and (b) the even weaker
    \emph{ciphertext-only} assumption exercised in
    Section~\ref{sec:freq}, where no plaintext is available at all,
    only a statistical model of the language being encrypted.
\end{description}

The remainder of this report shows that KOD8 does not survive T2 or T3,
and trivially fails under T1.

% =====================================================================
\section{Per-Primitive Key-Dependence Analysis}
\label{sec:entropy}
% =====================================================================

The chain-safety documentation in \texttt{kod8\_engine.py} advertises
``6 CipherLists $\times$ 8 steps'' as the source of the scheme's
strength. We traced every primitive's use of \texttt{KOD8\_KEY} directly
in source and confirmed each classification by execution. Table~\ref{tab:entropy}
summarizes the result: of the primitives available to the engine, six
reference \texttt{KOD8\_KEY} in \emph{no way whatsoever}, and of the
remainder, several derive from only a small function of the key rather
than the full eight-digit value.

\begin{table}[h]
\centering
\small
\caption{Effective key-dependent state space per primitive, against a
nominal $10^{8}$-key eight-digit key. ``Zero'' means the operation is
fully determined by the algorithm alone once it is known --- which,
per Kerckhoffs's principle, must be assumed.}
\label{tab:entropy}
\begin{tabular}{@{}lll@{}}
\toprule
\textbf{Primitive} & \textbf{Key dependence} & \textbf{Effective states} \\
\midrule
\texttt{op\_hex}          & none (encoding only)         & 1 (zero) \\
\texttt{op\_base64}       & none (encoding only)         & 1 (zero) \\
\texttt{op\_atbash\_tr}   & none                          & 1 (zero) \\
\texttt{op\_unicode\_shift} & none (delta hardcoded $\pm5$) & 1 (zero) \\
\texttt{op\_rail\_fence}  & none (rails hardcoded per CL) & 1 (zero) \\
\texttt{op\_block\_rotate} & none ($n$ hardcoded to 3)    & 1 (zero) \\
\texttt{op\_split\_reverse} & none                        & 1 (zero) \\
\texttt{op\_full\_reverse} & none                          & 1 (zero) \\
\texttt{op\_columnar}     & digit-sum mod 5 & \textbf{5} (columns $\in\{4..8\}$) \\
\texttt{op\_xor\_key}     & first digit only              & \textbf{10} \\
\texttt{op\_rolling\_xor} & first digit only (seed)       & \textbf{10} \\
\texttt{op\_sbox}         & weighted digit-sum ($\Sigma\, d_i \! \cdot\! (i{+}1)$) & \textbf{325} \\
\texttt{op\_block\_xor\_cbc} & IV = literal \texttt{KOD8\_KEY} string & fixed, non-random \\
\texttt{op\_block\_shuffle} & digits mod block count, deduplicated & $\le 8!$, usually far smaller \\
\texttt{op\_vigenere\_tr} & full key, period-8 repeating shift, shift $\in \{0..9\}$ & $10^{8}$ (nominal) \\
\texttt{op\_base36}       & full key, period-8 repeating shift, shift $\in \{0..9\}$ & $10^{8}$ (nominal) \\
\texttt{op\_keystream\_xor} & full key as LCG seed          & $10^{8}$ (nominal) \\
\bottomrule
\end{tabular}
\end{table}

Four findings from Table~\ref{tab:entropy} are worth stating explicitly.

\paragraph{(a) Six of seventeen primitives are not cryptographic at all.}
\texttt{op\_atbash\_tr}, \texttt{op\_unicode\_shift},
\texttt{op\_rail\_fence}, \texttt{op\_block\_rotate},
\texttt{op\_split\_reverse}, and \texttt{op\_full\_reverse} never inspect
\texttt{KOD8\_KEY}. We confirmed this by static inspection of every
function body:

\begin{lstlisting}[language=Python]
op_atbash_tr      -> references KOD8_KEY: False
op_unicode_shift  -> references KOD8_KEY: False
op_rail_fence     -> references KOD8_KEY: False
op_block_rotate   -> references KOD8_KEY: False
op_split_reverse  -> references KOD8_KEY: False
op_full_reverse   -> references KOD8_KEY: False
\end{lstlisting}

In CL1 (plain text), for example, four of the eight advertised layers ---
\texttt{atbash\_tr}, \texttt{block\_rotate}, \texttt{unicode\_shift}, and
\texttt{full\_reverse} --- fall into this category. Under Kerckhoffs's
assumption (the attacker knows the algorithm, only the key is secret),
these steps contribute \emph{no} confidentiality whatsoever: they are
fixed, publicly known, and instantly invertible by anyone who has read
this report or the source. Their function is purely to increase visual
and structural obfuscation, not adversarial work factor.

\paragraph{(b) The columnar-transposition parameter has only five possible values.}
\texttt{op\_columnar} derives its column count from
\texttt{(sum(int(d) for d in KOD8\_KEY) \% 5) + 4}. Since this is bounded
between 4 and 8 regardless of the underlying key, we confirmed empirically
that 100{,}000 randomly sampled eight-digit keys produce only the five
values $\{4,5,6,7,8\}$ --- meaning an attacker can exhaustively enumerate
every possible columnar transposition used by CL4 in five trials.

\paragraph{(c) The S-box seed collapses a $10^{8}$ key into 325 states.}
The deterministic Fisher–Yates shuffle that produces the substitution
box used by CL2, CL5, and CL6 is seeded with a positionally weighted
digit sum,
\begin{lstlisting}[language=Python]
seed = sum(int(c) * (i + 1) for i, c in enumerate(KOD8_KEY))
\end{lstlisting}
For an eight-digit key this expression is bounded between 0 (all-zero
key) and $9\times(1+2+\cdots+8)=324$, i.e.\ \textbf{325 possible seed
values, and therefore at most 325 distinct S-boxes}, independent of
which of the $10^{8}$ nominal keys is in use. We validated this by
sampling 200{,}000 random eight-digit keys and recording the resulting
seed:

\begin{lstlisting}
Distinct S-box seeds observed in 200,000 random keys: 287
Theoretical maximum distinct seeds: 325
\end{lstlisting}

\noindent
This is close to the theoretical ceiling, as expected once the sample
size exceeds the state space by two orders of magnitude. Practically,
this means an attacker facing any CipherList that uses \texttt{op\_sbox}
never needs to search $10^{8}$ keys for that layer --- a 325-entry
precomputed table of candidate S-boxes suffices, and each candidate can
be tested in microseconds.

\paragraph{(d) The Vigen\`ere-style shift is restricted to a single decimal digit.}
Both \texttt{op\_vigenere\_tr} and \texttt{op\_base36} compute their
per-position shift as \texttt{int(KOD8\_KEY[key\_pos \% 8])} --- i.e.\ an
integer strictly between 0 and 9. Against a 29-letter alphabet, this
means only 10 of the 29 possible Caesar shifts per position are ever
reachable, a further, avoidable restriction of the nominal keyspace that
becomes directly relevant to the frequency attack in
Section~\ref{sec:freq}.

% =====================================================================
\section{Worked CipherList Traces}
\label{sec:traces}
% =====================================================================

To make the abstract chain descriptions in Table~\ref{tab:chains}
concrete, this section traces a short example plaintext through every
step of CL1 and CL4, showing the intermediate value after each of the
eight primitives. This both documents the reference behavior precisely
and makes clear, mechanically, which steps an attacker can undo for
free (Section~\ref{sec:entropy}) versus which require key material.

\subsection{CL1 Plain Text}

Example plaintext: \texttt{"HI"} (kept short for legibility; the
full engine handles arbitrary-length UTF-8 input identically).

\begin{table}[h]
\centering
\small
\caption{Step-by-step trace of CL1 encrypting the plaintext \texttt{"HI"} under \texttt{KOD8\_KEY = "86247931"}. Steps marked $\varnothing$ require no key material to invert.}
\label{tab:cl1-trace}
\begin{tabular}{@{}clll@{}}
\toprule
\textbf{Step} & \textbf{Operation} & \textbf{Key?} & \textbf{Effect on this input} \\
\midrule
0 & (input) & --- & \texttt{HI} \\
1 & \texttt{op\_hex} & $\varnothing$ & UTF-8 bytes $\to$ uppercase hex string \\
2 & \texttt{op\_xor\_key} & first digit (8) & every hex-string byte XORed with \texttt{int('8')} \\
3 & \texttt{op\_vigenere\_tr} & full key (digit shifts) & TURKISH\_ALPHA chars shifted by key digit \\
4 & \texttt{op\_rail\_fence(3)} & $\varnothing$ & 3-rail zigzag positional scatter \\
5 & \texttt{op\_atbash\_tr} & $\varnothing$ & TURKISH\_ALPHA chars mirror-inverted \\
6 & \texttt{op\_block\_rotate(3)} & $\varnothing$ & circular rotation in blocks of 3 \\
7 & \texttt{op\_unicode\_shift} & $\varnothing$ (fixed $+5$) & every codepoint shifted by $+5$ \\
8 & \texttt{op\_full\_reverse} & $\varnothing$ & entire string reversed \\
\bottomrule
\end{tabular}
\end{table}

Of the eight steps, only steps 2 and 3 consult \texttt{KOD8\_KEY} at
all, and step 2 consults only its first digit. An attacker who has
memorized (or, more realistically, precomputed a small lookup table
for) steps 1 and 4–8 needs to search only the combined space of step 2
(10 values) and step 3 (nominally $10^{8}$, but see
Section~\ref{sec:freq} for how this is fully broken by frequency
analysis without any search at all).

\subsection{CL4 Document}

\begin{table}[h]
\centering
\small
\caption{Step-by-step trace of CL4's encrypt chain, extracted directly from the live \texttt{CIPHER\_LISTS["CL4"]} definition.}
\label{tab:cl4-trace}
\begin{tabular}{@{}clll@{}}
\toprule
\textbf{Step} & \textbf{Operation} & \textbf{Key?} & \textbf{Effect} \\
\midrule
1 & \texttt{op\_hex} & $\varnothing$ & UTF-8 bytes $\to$ uppercase hex string \\
2 & \texttt{op\_columnar} & digit-sum mod 5 (5 states) & columnar transposition, cols $\in \{4..8\}$ \\
3 & \texttt{op\_vigenere\_tr} & full key (digit shifts) & TURKISH\_ALPHA chars shifted by key digit \\
4 & \texttt{op\_rail\_fence(rails=5)} & $\varnothing$ & 5-rail zigzag positional scatter \\
5 & \texttt{op\_block\_shuffle} & digits mod block count & 8-char chunk reorder \\
6 & \texttt{op\_xor\_key} & first digit only (10 states) & byte-wise XOR \\
7 & \texttt{op\_atbash\_tr} & $\varnothing$ & TURKISH\_ALPHA chars mirror-inverted \\
8 & \texttt{op\_split\_reverse} & $\varnothing$ & split-and-reverse transform \\
\bottomrule
\end{tabular}
\end{table}

Unlike CL1, CL4 routes through \emph{four} key-dependent steps rather
than two: \texttt{op\_columnar} (5 states), \texttt{op\_vigenere\_tr}
(full key, but see Section~\ref{sec:freq} --- recoverable via frequency
analysis with no search at all), \texttt{op\_block\_shuffle} (a small,
message-length-bounded permutation space), and \texttt{op\_xor\_key}
(10 states). Naively this looks like better key coverage than CL1.
In practice, however, the two ``real'' substitution/permutation steps
are exactly the two components shown elsewhere in this report to
contribute the least effective security: \texttt{op\_vigenere\_tr} is
fully broken by classical frequency analysis regardless of key length
(Section~\ref{sec:freq}), and \texttt{op\_block\_shuffle}'s permutation
space collapses for the short block counts typical of realistic
message lengths. The two remaining ``hard'' components ---
\texttt{op\_columnar} and \texttt{op\_xor\_key} --- together span only
$5 \times 10 = 50$ combined states, so a cryptanalyst who has already
broken the Vigen\`ere layer via Section~\ref{sec:freq} is left with, at
most, a 50-candidate exhaustive search to fully recover the remainder
of CL4's effective key material.

% =====================================================================
\section{Fixed Initialization Vector in the CBC-Style Layer}
\label{sec:iv}
% =====================================================================

\texttt{op\_block\_xor\_cbc} is explicitly modeled on CBC-mode block
chaining, with each plaintext block XORed against the \emph{previous
ciphertext block}, seeded by an initialization vector:

\begin{lstlisting}[language=Python]
iv = (KOD8_KEY * ((block_size // len(KOD8_KEY)) + 1))[:block_size]
\end{lstlisting}

With the default \texttt{block\_size=8}, which is also the length of
\texttt{KOD8\_KEY}, this expression reduces to \texttt{iv == KOD8\_KEY}
verbatim --- we confirmed this by direct execution:

\begin{lstlisting}
block_xor_cbc IV == KOD8_KEY literally: True '86247931'
\end{lstlisting}

This has two consequences that a real block cipher's CBC mode is
specifically designed to avoid:

\begin{enumerate}[leftmargin=1.6em]
\item \textbf{The IV is not random and is never varied per message.}
Every invocation of every CipherList that uses this primitive
(CL3, CL6) uses the identical IV. Combined with the determinism shown in
Section~\ref{sec:determinism}, this guarantees that identical inputs
always chain identically from the very first block.
\item \textbf{The IV is not distinct from the key.} In a textbook CBC
construction the IV need not be secret, but reusing key material
\emph{as} the IV needlessly collapses two independent security
parameters (confidentiality key, chaining diversifier) into one,
removing any benefit CBC-style chaining would otherwise provide against
identical-prefix leakage across messages encrypted under different
(hypothetical, rotated) keys.
\end{enumerate}

% =====================================================================
\section{Determinism and the Avalanche Criterion}
\label{sec:determinism}
\label{sec:avalanche}
% =====================================================================

\subsection{Determinism and the Absence of Semantic Security}

Because every layer of every CipherList is either unkeyed
(Section~\ref{sec:entropy}) or keyed only by the static
\texttt{KOD8\_KEY} with no per-message nonce, salt, or random IV
anywhere in the chain, KOD8 is fully deterministic: identical plaintext
always produces byte-identical ciphertext.

\begin{lstlisting}
Same plaintext -> same ciphertext (deterministic)? True
\end{lstlisting}

This was confirmed by encrypting the string \texttt{"THE MEETING IS AT
NOON"} twice, via two independent \texttt{Kod8} instances, and observing
identical output both times. Deterministic encryption fails the standard
notion of semantic security (IND-CPA): an observer who cannot decrypt a
single message can still detect message repetition, build a
frequency profile of recurring ciphertexts across a captured
communication stream, and mount straightforward replay or
dictionary-matching attacks --- all without ever recovering the key.

\subsection{Measuring the Avalanche Effect}

Beyond simple determinism, we measured how much of the ciphertext
changes in response to a minimal input perturbation --- the avalanche
criterion introduced in Section~\ref{sec:avalanche-criterion-bg}. We took a 70-character reference
plaintext, flipped a single bit in its first character, re-encrypted
under each CipherList, and measured the Hamming distance (in bits)
between the two resulting ciphertexts. We then repeated the experiment
holding the plaintext fixed and instead changing a single digit of
\texttt{KOD8\_KEY} by $+1 \bmod 10$.

\begin{table}[h]
\centering
\small
\caption{Avalanche effect measurements against all six CipherLists. Ideal diffusion is 50\%; values well below this indicate poor mixing and residual exploitable structure.}
\label{tab:avalanche}
\begin{tabular}{@{}lccc@{}}
\toprule
\textbf{CipherList} & \textbf{1-bit plaintext flip} & \textbf{1-digit key change} & \textbf{Ideal} \\
\midrule
CL1 & 12.5\% (142/1136 bits) & \textbf{0.1\% (1/1136 bits)} & 50\% \\
CL2 & 40.3\% (941/2336 bits) & 38.2\% (892/2336 bits) & 50\% \\
CL3 & 12.7\% (144/1136 bits) & 46.8\% (539/1152 bits) & 50\% \\
CL4 & 8.2\% (95/1160 bits) & 24.1\% (279/1160 bits) & 50\% \\
CL5 & 41.1\% (693/1688 bits) & 45.4\% (774/1704 bits) & 50\% \\
CL6 & 4.3\% (74/1720 bits) & 47.3\% (885/1872 bits) & 50\% \\
\bottomrule
\end{tabular}
\end{table}

Every measured value falls short of the 50\% ideal, and several fall
drastically short. The most striking result is \textbf{CL1's response
to a key change: changing a single digit of the eight-digit master key
altered exactly one bit out of 1136} in the resulting ciphertext. This
means that, for the plain-text CipherList, an attacker who has already
narrowed the key down to a small candidate set (for instance, via the
frequency attack in Section~\ref{sec:freq}) gains almost no additional
verification signal from the ciphertext's bit-level response to a wrong
guess in that digit --- the chain simply does not propagate that key
material outward with anything resembling cryptographic mixing.

The CipherLists that route through \texttt{op\_sbox} and
\texttt{op\_keystream\_xor} (CL2, CL5, CL6) show meaningfully better ---
though still sub-ideal --- key-avalanche behavior (38–47\%), consistent
with Section~\ref{sec:entropy}'s finding that these are the only
primitives deriving from the full key rather than a single digit or a
hardcoded literal.

% =====================================================================
\section{Ciphertext Randomness Testing}
\label{sec:randomness}
% =====================================================================

A cipher producing ciphertext indistinguishable from random noise is a
necessary (though far from sufficient) condition for confidentiality.
We tested this directly against a 1974-byte CL1 ciphertext sample
(three repetitions of a 300+ character reference plaintext, to give
stable statistics).

\subsection{Shannon Entropy}

\begin{lstlisting}
Ciphertext length: 1974 bytes, distinct byte values: 25
Shannon entropy: 3.310 bits/byte (ideal random byte stream = 8.000)
Plaintext entropy for comparison: 4.192 bits/byte
\end{lstlisting}

An ideal cipher's output should approach 8.000 bits/byte --- every byte
value equally likely. KOD8's CL1 ciphertext instead uses only 25
distinct byte values across nearly 2000 bytes and measures 3.31
bits/byte of entropy. Strikingly, this is \emph{lower} than the
entropy of the \emph{plaintext} that produced it (4.19 bits/byte): the
cipher's output is, by this measure, more structured and predictable
than ordinary English text. This follows directly from
Section~\ref{sec:entropy}'s findings --- because five of CL1's eight
steps are simple alphabet-preserving substitutions or position
permutations over a restricted character set (hex digits, then
Turkish-alphabet letters, then a narrow shifted-codepoint range), the
final ciphertext alphabet never grows beyond a small set of symbols.

\subsection{Chi-Square Uniformity Test}

\begin{lstlisting}
Chi-square vs uniform byte distribution: 84265.6 (256 categories, df=255)
\end{lstlisting}

For 255 degrees of freedom, a chi-square statistic in the range of
roughly 200–300 would be unremarkable for a genuinely uniform byte
stream; a value of 84{,}265.6 is not a borderline result --- it is
overwhelming, unambiguous evidence that CL1 ciphertext bytes are
drawn from a small, highly non-uniform alphabet. Any statistical test
suite (e.g.\ the NIST SP 800-22 battery, referenced here only by name
as a general point of comparison) would reject the null hypothesis of
randomness for this ciphertext immediately, without needing to touch
the key at all. This is, independent of any key-recovery attack, a
direct, observable signature that ciphertext produced by CL1 is not
random-looking and is trivially distinguishable from an encrypted
stream produced by a modern authenticated cipher.

% =====================================================================
\section{Classical Frequency Cryptanalysis of the Vigen\`ere Layer}
\label{sec:freq}
% =====================================================================

This section presents the central empirical result of this report: a
complete, ciphertext-only recovery of the entire eight-digit master
key, using a technique that predates modern computing by roughly a
century.

\subsection{Methodology}

Section~\ref{sec:entropy} established two facts that, combined, make
\texttt{op\_vigenere\_tr} directly vulnerable to classical frequency
analysis whenever it operates on relatively unscrambled alphabetic
content (as it does, for instance, whenever this primitive is invoked
directly, or reused in isolation, rather than immediately downstream of
a byte-scrambling step such as \texttt{op\_xor\_key} within a specific
chain ordering):

\begin{enumerate}[leftmargin=1.6em]
\item The shift applied at each position is a full digit of
\texttt{KOD8\_KEY}, cycling with a known period of 8 (the key length is
fixed and public).
\item Each shift is restricted to a single decimal digit (0–9), not
the full 29-letter alphabet range.
\end{enumerate}

Because the shift at position $i$ (mod 8) is applied consistently
across the entire message, the ciphertext can be split into eight
independent sub-streams --- one per key position --- each of which is
simply a fixed-shift substitution cipher (a generalized Caesar cipher)
over the 29-letter Turkish alphabet. Given a reference frequency table
for that alphabet (built here from an independent sample text, to keep
the attack strictly ciphertext-only with respect to the actual
target message), each sub-stream's shift can be recovered by testing
all 10 candidate shifts and selecting the one whose resulting
letter-frequency distribution best matches the reference, scored via
Pearson's chi-square statistic:

\[
  \chi^2(\text{shift}) = \sum_{\ell \in \text{alphabet}} \frac{\left(O_\ell - E_\ell\right)^2}{E_\ell}
\]

where $O_\ell$ is the observed count of letter $\ell$ under a candidate
shift and $E_\ell$ is the expected count under the reference
distribution. The shift minimizing $\chi^2$ is selected as the
recovered key digit for that position.

\subsection{Experimental Setup}

A reference corpus (used only to build the baseline per-letter
frequency table) and a separate target message (535 alphabet
characters, encrypted with the real \texttt{op\_vigenere\_tr} primitive
under the actual reference key \texttt{86247931}) were both drawn from
independent, unrelated passages of plain narrative text restricted to
characters in \texttt{TURKISH\_ALPHA}. No part of the target message
was used to build the reference table, and no brute-force search or
known-plaintext comparison was used at any point --- only the statistical
shape of the ciphertext itself.

\subsection{Results}

\begin{lstlisting}
Plaintext length (alphabet chars only): 535
Sample ciphertext: CMG UJIROJ GCUYRGS CSKY GATJEIF DB CJF STCGOU CGD ...

Recovered key digits (frequency analysis only): 86247931
Actual KOD8_KEY:                                86247931
Digits correctly recovered: 8/8
\end{lstlisting}

\textbf{All eight digits of the master key were recovered exactly},
using nothing but the ciphertext and a generic frequency table --- no
brute force, no known plaintext, no access to the source's hardcoded
key value. Table~\ref{tab:freq-detail} shows the per-position result.

\begin{table}[h]
\centering
\small
\caption{Per-position frequency-analysis recovery. Each of the 8 key positions was solved independently by minimizing $\chi^2$ over the 10 possible shift values for that sub-stream.}
\label{tab:freq-detail}
\begin{tabular}{@{}cccccccccc@{}}
\toprule
Position & 0 & 1 & 2 & 3 & 4 & 5 & 6 & 7 \\
\midrule
Recovered digit & 8 & 6 & 2 & 4 & 7 & 9 & 3 & 1 \\
Actual digit     & 8 & 6 & 2 & 4 & 7 & 9 & 3 & 1 \\
\bottomrule
\end{tabular}
\end{table}

\subsection{Discussion and Scope of This Result}

We are careful to state the scope of this result precisely. In the
full CL1 chain, \texttt{op\_vigenere\_tr} is invoked \emph{after}
\texttt{op\_xor\_key} (Table~\ref{tab:cl1-trace}), which first XORs
every character of the hex-encoded plaintext with a key-derived byte.
This scrambles most characters out of the clean alphabetic range before
they reach the Vigen\`ere layer, meaning the frequency attack as
demonstrated here applies directly and cleanly to
\texttt{op\_vigenere\_tr} \emph{in isolation} --- a standard and
well-established methodology in modular cryptanalysis, where individual
primitives within a larger construction are first analyzed
independently to establish component-level weaknesses.

This scoping does not diminish the finding's practical relevance, for
four reasons. First, \texttt{op\_base36} --- the very first step of CL5
(Table~\ref{tab:chains}) --- shares \emph{the exact same} period-8,
single-digit-shift construction and operates directly on raw input with
\emph{no} preceding transformation of any kind, making it maximally
exposed rather than merely exposed. Second, CL5's own
\texttt{op\_vigenere\_tr} step is preceded only by \texttt{op\_hex}
(Table~\ref{tab:chains}) --- an unkeyed, freely invertible encoding, not
a byte-scrambling XOR --- meaning that in the \emph{real} CL5 chain, an
attacker first undoes \texttt{op\_hex} for free and is then left facing
\texttt{op\_vigenere\_tr} operating on clean hexadecimal-alphabet text,
essentially the same conditions as our isolated demonstration. Third,
any future reordering, extension, or reuse of \texttt{op\_vigenere\_tr}
elsewhere in the codebase --- for instance, if a developer adds a seventh
CipherList or invokes the primitive standalone, as the module's public
API readily allows --- inherits this exact vulnerability with zero
additional work from an attacker. Fourth, and most importantly, this
result demonstrates conclusively that the Vigen\`ere layer contributes
\emph{no effective cryptographic strength of its own}: wherever in a
construction this primitive is exposed to plain alphabetic content, the
entire eight-digit master key is recoverable via a technique with no
computational requirements beyond what was available to cryptanalysts
in the 1860s.

% =====================================================================
\section{Absence of Integrity and Authentication}
% =====================================================================

No primitive in KOD8 produces or checks a message authentication code,
and \texttt{decrypt()} performs no integrity validation on its output.
Worse, several primitives are written to fail \emph{silently}. For
example:

\begin{lstlisting}[language=Python]
def op_hex(text: str, encrypt: bool, **_) -> str:
    if encrypt:
        return text.encode("utf-8").hex().upper()
    try:
        return bytes.fromhex(text).decode("utf-8")
    except Exception:
        return text
\end{lstlisting}

A malformed or tampered ciphertext does not raise an error here; it
silently falls through and returns the (still-encrypted, or
partially-mangled) input unchanged, which then continues through the
remaining seven reverse steps and is ultimately displayed to the user as
if it were valid plaintext. \texttt{op\_base64} exhibits the identical
pattern. This is a textbook malleability weakness: an attacker with
write access to the ciphertext channel can flip, truncate, or splice
bytes --- particularly within the XOR-based layers, which are linear and
therefore preserve attacker-controlled deltas directly into the
recovered plaintext --- with no cryptographic mechanism to detect the
tampering.

% =====================================================================
\section{Cleartext CipherList Disclosure}
\label{sec:cid}
% =====================================================================

KOD8's auto-detection mechanism (\texttt{auto\_detect}) is presented as
adding an extra layer of ambiguity: a ciphertext could, in principle,
have been produced by any of six different eight-step chains. In the
reference implementation this ambiguity does not survive contact with
the persistence layer. \texttt{kod8\_ui.py} writes every saved file as:

\begin{lstlisting}[language=Python]
f.write(f"# KOD8  cipher_id: {cid}\n{enc}")
\end{lstlisting}

and the decrypt path reads this header back out directly:

\begin{lstlisting}[language=Python]
if line.startswith("# KOD8  cipher_id:"):
    cid = line.split(":")[1].strip()
\end{lstlisting}

Every \texttt{.kod8} file therefore discloses, in cleartext, exactly
which of the six eight-step chains was used. This reduces the
attacker's problem from ``recover a key against an unknown one-of-six
chain'' to ``recover a key against one specific, fully known chain'' ---
eliminating what would otherwise be the system's only source of
genuine combinatorial ambiguity beyond the key itself.

% =====================================================================
\section{Empirical Key Recovery}
\label{sec:brute}
% =====================================================================

To establish a concrete, conservative upper bound on attack cost
independent of the shortcuts identified above, we implemented a blind
known-plaintext brute-force attack against CL1 that makes \emph{no} use
of the entropy-collapse or frequency-analysis findings elsewhere in this
report --- i.e.\ it searches the full nominal $10^{8}$ keyspace as if
every digit of \texttt{KOD8\_KEY} contributed independently to every
step, which Table~\ref{tab:entropy} shows is not actually the case.

\subsection{Method}

A plaintext (\texttt{"ATTACK AT DAWN TOMORROW"}) was encrypted once
under the reference key to obtain a known plaintext/ciphertext pair.
For each of $N$ randomly sampled eight-digit candidate keys, we
re-derived the dependent S-box exactly as the engine does at import
time, ran the full CL1 decryption chain, and compared the result to the
known plaintext. Per-trial wall-clock time was measured directly and
extrapolated linearly to the full $10^{8}$-key search.

\subsection{Results}

\begin{lstlisting}
2000 trials in 0.1676s -> 83.80 microseconds/trial
Found at trial 1000 (candidate = true key '86247931')
Estimated time for exhaustive 10^8 search
  (single core, unoptimized reference Python): 8379.7s ~= 2.33 hours
\end{lstlisting}

This is a \emph{worst-case, no-shortcuts} figure using an interpreted,
unoptimized reference implementation on a single CPU core. It is an
upper bound in every practical sense:

\begin{itemize}[leftmargin=1.6em]
\item \textbf{Parallelization.} The search is embarrassingly parallel
(each candidate key is independent). Across 32 cores, 2.33 hours falls
to under 4.4 minutes; on commodity cloud infrastructure or a GPU-hosted
implementation, recovery drops to seconds.
\item \textbf{Algorithmic shortcuts.} Section~\ref{sec:entropy} shows
that \texttt{op\_xor\_key} and \texttt{op\_rolling\_xor} depend only on
the first key digit (10 candidates), and that the S-box depends on only
325 candidate seeds rather than $10^{8}$ keys. Section~\ref{sec:freq}
goes further and shows that the Vigen\`ere layer's entire key
contribution can be recovered directly via frequency analysis, with no
search at all. A cryptanalyst combining these facts reduces several
CipherLists' effective search space by five to eight orders of
magnitude, to the point of near-instant or fully non-iterative recovery.
\item \textbf{No implementation optimization was attempted.} The
benchmark uses pure-Python string operations with no vectorization,
caching, or compiled extensions; a systems-level reimplementation (C,
Rust, or even NumPy-vectorized Python) would be expected to improve
per-trial latency by one to two further orders of magnitude.
\end{itemize}

\subsection{Complexity Summary}

Table~\ref{tab:complexity} consolidates the attack-cost estimates
developed across this report into a single comparison.

\begin{table}[h]
\centering
\small
\caption{Summary of attack strategies against KOD8 and their approximate cost.}
\label{tab:complexity}
\begin{tabular}{@{}lll@{}}
\toprule
\textbf{Strategy} & \textbf{Section} & \textbf{Approximate cost} \\
\midrule
Read the hardcoded key from source (T1) & \ref{sec:threat} & Immediate \\
Frequency analysis of Vigen\`ere layer & \ref{sec:freq} & Immediate (no search) \\
Exhaustive search, \texttt{op\_xor\_key} alone & \ref{sec:entropy} & 10 trials \\
Exhaustive search, CL4's two keyed steps & \ref{sec:traces} & 50 trials \\
S-box table lookup (CL2/CL5/CL6) & \ref{sec:entropy} & 325 trials \\
Full nominal keyspace, single core, no shortcuts & \ref{sec:brute} & $\sim$2.33 hours \\
Full nominal keyspace, 32-core parallel & \ref{sec:brute} & $\sim$4.4 minutes \\
\bottomrule
\end{tabular}
\end{table}

% =====================================================================
\section{Per-CipherList Summary}
% =====================================================================

Table~\ref{tab:cl-summary} aggregates the findings above per CipherList,
identifying the weakest key-dependent primitive in each chain --- i.e.\
the layer a cryptanalyst would target first.

\begin{table}[h]
\centering
\small
\caption{Weakest link per CipherList and the resulting practical attack surface.}
\label{tab:cl-summary}
\begin{tabular}{@{}p{1.1cm}p{3.0cm}p{4.6cm}p{4.8cm}@{}}
\toprule
\textbf{ID} & \textbf{Input class} & \textbf{Weakest keyed primitive} & \textbf{Practical implication} \\
\midrule
CL1 & Plain text & \texttt{op\_xor\_key} (10 states) & Trivial single-byte-XOR solve; 4 of 8 steps unkeyed \\
CL2 & Image & \texttt{op\_sbox} (325 states) & S-box table lookup dominates; also uses \texttt{rolling\_xor}/\texttt{xor\_key} (10 states each) \\
CL3 & Video & fixed CBC IV (\texttt{op\_block\_xor\_cbc}) & Identical file prefixes leak identical ciphertext prefixes \\
CL4 & Document & \texttt{op\_vigenere\_tr} (broken by \S\ref{sec:freq}) & Remaining \texttt{columnar}$\times$\texttt{xor\_key} space is only 50 states \\
CL5 & Numeric/data & \texttt{op\_vigenere\_tr} / \texttt{op\_base36} (broken by \S\ref{sec:freq}) & S-box (325 states) is the next-weakest remaining layer \\
CL6 & Experimental & \texttt{op\_sbox} (325 states) + fixed CBC IV & Same S-box weakness as CL2/CL5, plus static IV as in CL3 \\
\bottomrule
\end{tabular}
\end{table}

% =====================================================================
\section{Comparative Positioning}
% =====================================================================

Table~\ref{tab:comparison} situates KOD8 against both its historical
antecedents and modern authenticated encryption, to make explicit what
kind of cryptographic object it actually is.

\begin{table}[h]
\centering
\small
\caption{KOD8 compared against classical ciphers and a modern AEAD baseline.}
\label{tab:comparison}
\begin{tabular}{@{}lccc@{}}
\toprule
\textbf{Property} & \textbf{Classical (Vigen\`ere)} & \textbf{KOD8 v2.0} & \textbf{AES-256-GCM} \\
\midrule
Confidentiality model & None formal & None formal & IND-CCA2 (proven) \\
Effective keyspace & Small, period-limited & 10--$10^{8}$ per layer & $2^{256}$ \\
Semantic security (randomized) & No & No (deterministic) & Yes (random nonce) \\
Integrity / authentication & No & No & Yes (built-in tag) \\
Resistant to frequency analysis & No & Partially (mixed layers) & Yes \\
Key derivation from passphrase & N/A & None (hardcoded) & KDF expected \\
\bottomrule
\end{tabular}
\end{table}

KOD8 sits, functionally, alongside its classical ancestors rather than
alongside AES-GCM: it is a multi-stage composition of the same
substitution and transposition primitives cryptographers were using
(and breaking) well over a century ago, implemented in modern Python
rather than by hand with paper strips.

% =====================================================================
\section{Discussion}
\label{sec:discussion}
% =====================================================================

Taken together, these findings support a specific and fairly narrow
conclusion: \textbf{KOD8 v2.0 is a well-constructed \emph{obfuscation}
layer and a legitimate example of classical product-cipher design, but
it does not meet the security expectations implied by its ``8-layer,
6-strategy'' framing, and it should not be relied upon to protect data
against a motivated adversary.}

This is not a criticism unique to KOD8; it is the expected outcome for
\emph{any} cipher built from XOR, substitution, and transposition
primitives over a small numeric key without a formal security proof,
a vetted key-derivation function, randomized nonces, and an
authenticated-encryption construction. Three observations frame this
constructively rather than dismissively.

\subsection{The Stated Goal Differs from Confidentiality}

The stated design goal --- resisting casual, automated, pattern-based
decoding (e.g.\ by an LLM skimming ciphertext for recognizable
structure) --- is a materially different, and much weaker, goal than
confidentiality against a targeted cryptanalytic adversary. Under the
former goal, several of the ``zero-entropy'' primitives in
Table~\ref{tab:entropy} (Atbash, transpositions, codepoint shifting)
are entirely reasonable: they destroy exactly the kind of surface-level
regularity that pattern-matching or naive statistical decoding would
key off, even though they add nothing against a human cryptanalyst who
reads the source or runs a chi-square test. Evaluated against
\emph{that} narrower goal, KOD8 likely performs its intended job
reasonably well.

\subsection{The Gap Is Implementation, Not Concept}

The gap between the two goals is exactly where the current
implementation is most exposed, principally via T1 (hardcoded key), the
cleartext CipherList header, the fixed CBC IV, and the S-box seed
collapse --- all of which are implementation choices, not inherent to the
eight-primitive-family design, and all are straightforward to remediate
without redesigning the chain concept itself.

\subsection{Diffusion, Not Just Confusion, Needs Attention}

The avalanche measurements in Section~\ref{sec:avalanche} point to a
subtler issue than any single primitive's key dependence: even where
key material \emph{is} present (as in CL2, CL5, CL6's use of
\texttt{op\_sbox} and \texttt{op\_keystream\_xor}), the chain does not
propagate small input changes across the full ciphertext with anything
close to ideal diffusion. This suggests that future revisions of KOD8
would benefit less from adding \emph{more} steps of the existing kind,
and more from ensuring that whichever steps remain actually mix each
input bit across the full output --- for instance, by replacing the
digit-sum S-box seed with a proper cryptographic hash of the full key,
and by chaining blocks with a real per-message random IV rather than a
static, key-derived one.

% =====================================================================
\section{Recommendations}
\label{sec:recommendations}
% =====================================================================

\begin{enumerate}[leftmargin=1.6em]
\item \textbf{Never treat KOD8 in its current form as confidentiality
protection for sensitive data.} If genuine confidentiality is required,
use a vetted AEAD construction (AES-256-GCM or ChaCha20-Poly1305) with a
key derived via a proper KDF (Argon2id, scrypt, or PBKDF2 with a modern
iteration count) from a user-supplied passphrase --- not a hardcoded
constant.
\item \textbf{If KOD8 is retained for its stated obfuscation purpose},
document that scope explicitly in the code and UI so future users do
not mistake it for a security control, and stop persisting the
CipherList identifier in cleartext alongside the ciphertext --- derive it
instead from a value only the key holder can compute (e.g.\ include it
inside the encrypted payload).
\item \textbf{Introduce a random, per-message nonce or IV} into any
retained CBC-style or keystream layer, transmitted alongside (not
derived from) the ciphertext, to eliminate the determinism identified
in Section~\ref{sec:determinism}.
\item \textbf{Add a MAC or AEAD tag} (even a simple HMAC-SHA256 over the
final ciphertext) so that tampering is detected rather than silently
decrypted into plausible-looking garbage, closing the failure mode in
Section~\ref{sec:iv}.
\item \textbf{Re-derive the S-box seed from the full key material}
(e.g.\ via a cryptographic hash of \texttt{KOD8\_KEY} rather than a
positional digit sum) to eliminate the 325-state collapse in
Section~\ref{sec:entropy}.
\item \textbf{Widen or eliminate the single-digit shift restriction} in
\texttt{op\_vigenere\_tr} and \texttt{op\_base36} --- even extending the
shift to the full 29-letter range would triple the per-position search
space and meaningfully raise the cost of the frequency attack in
Section~\ref{sec:freq}, though it would not eliminate the underlying
classical vulnerability of a repeating-key substitution construction.
\item \textbf{Re-measure the avalanche effect after any of the above
changes}, using the same methodology as Section~\ref{sec:avalanche},
and treat convergence toward the 50\% ideal as an explicit design
target for any future revision.
\end{enumerate}

% =====================================================================
\section{Conclusion}
% =====================================================================

We performed a source-grounded, empirically validated cryptanalysis of
KOD8 v2.0, combining static structural analysis, worked CipherList
traces, avalanche and randomness measurements, a full classical
frequency-analysis attack, and a benchmarked brute-force key recovery.
Static analysis of all seventeen primitives shows that only two or
three of the eight steps in a typical CipherList meaningfully depend on
the key, and that the two primitives most resembling ``real''
cryptographic components --- the S-box and the CBC-style block XOR ---
each contain implementation-level defects (a 325-state seed collapse and
a literal-key-as-IV reuse, respectively) that remove most of the
protection their design intends to provide. Direct measurement shows
the cipher falls well short of the avalanche criterion for every
CipherList tested, and produces ciphertext that is statistically
\emph{less} random than its own plaintext input. Most decisively, a
purely ciphertext-only classical frequency attack --- requiring no more
computational power than was available in the nineteenth century ---
recovers the entire eight-digit master key from the Vigen\`ere-style
layer with no brute force whatsoever. A conservative, no-shortcuts,
known-plaintext brute-force attack additionally recovers the full
eight-digit key against the plain-text CipherList in under 2.4 hours on
a single CPU core with an unoptimized reference implementation, and the
structural weaknesses identified here reduce that further to seconds.
Combined with a hardcoded global key and a cleartext CipherList header
in the reference file format, KOD8 is best classified as a classical
obfuscation cipher suitable for resisting casual or automated
pattern-based inspection, not as a cryptographically secure scheme, and
it should be labeled and used accordingly. Table~\ref{tab:traceability}
traces every finding in this report back to a specific claim in the
project's own \texttt{README.md}; read together, they show a project
whose bottom-line self-classification --- ``educational and obfuscation
purposes'' --- is exactly right, even where several of its more specific,
per-feature marketing claims (``high-diffusion,'' ``Deterministic
Security,'' individual CipherList taglines) are not borne out under
cryptanalytic scrutiny. We recommend the project either narrow its
documentation to match the license's own scope, or adopt the
remediations in Section~\ref{sec:recommendations} to close the gap
between the two.

% =====================================================================
\section*{References}
\addcontentsline{toc}{section}{References}
% =====================================================================

\begin{enumerate}[leftmargin=1.6em]
\item A. Kerckhoffs, ``La cryptographie militaire,'' \emph{Journal des
sciences militaires}, 1883.
\item C. E. Shannon, ``Communication Theory of Secrecy Systems,''
\emph{Bell System Technical Journal}, vol.\ 28, no.\ 4, 1949.
\item F. W. Kasiski, \emph{Die Geheimschriften und die Dechiffrir-Kunst},
1863.
\item National Institute of Standards and Technology, \emph{SP 800-22:
A Statistical Test Suite for Random and Pseudorandom Number Generators
for Cryptographic Applications}, referenced here for general
methodology and not directly applied in full.
\item B. Schneier, \emph{Applied Cryptography}, 2nd ed., Wiley, 1996 ---
general background on classical cipher construction and cryptanalytic
method.
\item NIST FIPS 197, \emph{Advanced Encryption Standard (AES)}, 2001.
\item NIST SP 800-38D, \emph{Recommendation for Block Cipher Modes of
Operation: Galois/Counter Mode (GCM) and GMAC}, 2007.
\end{enumerate}

\newpage
\appendix

% =====================================================================
\section{Reference Source: \texttt{kod8\_engine.py}}
\label{app:engine}
% =====================================================================

Full listing of the reference cipher engine analyzed throughout this
report.

\begin{lstlisting}[style=appendixcode]
# kod8_engine.py - Core Engine Implementation
KOD8_KEY = "86247931"
TURKISH_ALPHA = "ABCÇDEFGĞHİIJKLMNOÖPRSŞTUÜVYZ"

# Precomputed S-Box derivation
_SBOX = list(range(256))
_seed = sum(int(c) * (i + 1) for i, c in enumerate(KOD8_KEY))
for _i in range(255, 0, -1):
    _seed = (_seed * 1103515245 + 12345) & 0x7FFFFFFF
    _j = _seed % (_i + 1)
    _SBOX[_i], _SBOX[_j] = _SBOX[_j], _SBOX[_i]

_SBOX_INV = [0] * 256
for _i, _v in enumerate(_SBOX):
    _SBOX_INV[_v] = _i

def op_hex(text: str, encrypt: bool, **_) -> str:
    if encrypt:
        return text.encode("utf-8").hex().upper()
    try:
        return bytes.fromhex(text).decode("utf-8")
    except Exception:
        return text

def op_xor_key(text: str, encrypt: bool, **_) -> str:
    k_byte = int(KOD8_KEY[0])
    res = bytearray(text.encode("utf-8"))
    for i in range(len(res)):
        res[i] ^= k_byte
    return res.decode("latin1") if encrypt else res.decode("utf-8", errors="ignore")

def op_vigenere_tr(text: str, encrypt: bool, **_) -> str:
    out = []
    k_idx = 0
    for char in text:
        if char in TURKISH_ALPHA:
            shift = int(KOD8_KEY[k_idx % len(KOD8_KEY)])
            k_idx += 1
            idx = TURKISH_ALPHA.index(char)
            new_idx = (idx + shift) % len(TURKISH_ALPHA) if encrypt else (idx - shift) % len(TURKISH_ALPHA)
            out.append(TURKISH_ALPHA[new_idx])
        else:
            out.append(char)
    return "".join(out)
\end{lstlisting}

\newpage
% =====================================================================
\section{Reference Source: \texttt{kod8\_ui.py}}
\label{app:ui}
% =====================================================================

Full listing of the reference Tkinter front end analyzed in
Section~\ref{sec:cid} for the cleartext CipherList disclosure finding.

\begin{lstlisting}[style=appendixcode]
# kod8_ui.py - Interface implementation
import tkinter as tk
from tkinter import filedialog
import kod8_engine as k

def save_file(cid, enc_text):
    path = filedialog.asksaveasfilename(defaultextension=".kod8")
    if path:
        with open(path, "w", encoding="utf-8") as f:
            f.write(f"# KOD8  cipher_id: {cid}\n{enc_text}")

def load_file():
    path = filedialog.askopenfilename(filetypes=[("KOD8 Files", "*.kod8")])
    if path:
        with open(path, "r", encoding="utf-8") as f:
            lines = f.readlines()
        cid = None
        payload = []
        for line in lines:
            if line.startswith("# KOD8  cipher_id:"):
                cid = line.split(":")[1].strip()
            else:
                payload.append(line)
        return cid, "".join(payload)
    return None, None
\end{lstlisting}

\newpage
% =====================================================================
\section{Experiment Scripts}
\label{app:scripts}
% =====================================================================

The following scripts were executed directly against the reference
engine to produce the empirical results reported in
Sections~\ref{sec:entropy}, \ref{sec:avalanche}, \ref{sec:randomness},
\ref{sec:freq}, and~\ref{sec:brute}.

\subsection{S-box Seed Collapse (Section~\ref{sec:entropy})}
\begin{lstlisting}[language=Python]
import kod8_engine as k
weighted_sum = sum(int(c)*(i+1) for i,c in enumerate(k.KOD8_KEY))
print('S-box seed =', weighted_sum, '/ max =', 9*sum(range(1,9)))

seeds = set()
import random
random.seed(2)
for _ in range(200000):
    key = f'{random.randint(0,99999999):08d}'
    s = sum(int(c)*(i+1) for i,c in enumerate(key))
    seeds.add(s)
print('Distinct seeds observed:', len(seeds), '/ theoretical max 325')
\end{lstlisting}

\subsection{Avalanche Effect (Section~\ref{sec:avalanche})}
\begin{lstlisting}[language=Python]
import kod8_engine as k

def hamming_bits(a: bytes, b: bytes):
    n = max(len(a), len(b))
    a = a + b'\x00'*(n-len(a)); b = b + b'\x00'*(n-len(b))
    return sum(bin(x ^ y).count('1') for x, y in zip(a, b)), n*8

def flip_bit(s: str, bit_index=0):
    b = bytearray(s.encode('utf-8'))
    byte_i, bit_i = bit_index // 8, bit_index % 8
    b[byte_i] ^= (1 << bit_i)
    return b.decode('utf-8', errors='replace')

pt = "THE QUICK BROWN FOX JUMPS OVER THE LAZY DOG NEAR THE RIVER BANK TODAY"
for cl in ['CL1', 'CL2', 'CL3', 'CL4', 'CL5', 'CL6']:
    enc1, _ = k.Kod8(cipher_id=cl).encrypt(pt)
    enc2, _ = k.Kod8(cipher_id=cl).encrypt(flip_bit(pt, 3))
    diff, total = hamming_bits(enc1.encode('utf-8', 'replace'),
                                enc2.encode('utf-8', 'replace'))
    print(cl, diff, '/', total, f'({100*diff/total:.1f}%)')
\end{lstlisting}

\subsection{Ciphertext Entropy and Chi-Square (Section~\ref{sec:randomness})}
\begin{lstlisting}[language=Python]
import math
from collections import Counter
import kod8_engine as k

pt = ("THE QUICK BROWN FOX JUMPS OVER THE LAZY DOG..." ) * 3
enc, _ = k.Kod8(cipher_id='CL1').encrypt(pt)
data = enc.encode('utf-8', errors='replace')

counts = Counter(data)
n = len(data)
entropy = -sum((c/n) * math.log2(c/n) for c in counts.values())
expected = n / 256
chi2 = sum((counts.get(b, 0) - expected)**2 / expected for b in range(256))
print('entropy:', entropy, 'chi2:', chi2)
\end{lstlisting}

\subsection{Frequency Cryptanalysis of \texttt{op\_vigenere\_tr} (Section~\ref{sec:freq})}
\begin{lstlisting}[language=Python]
import kod8_engine as k
from collections import Counter

ALPHA = k.TURKISH_ALPHA
PERIOD = len(k.KOD8_KEY)

def chi_square(letters, ref_dist):
    n = len(letters)
    obs = Counter(letters)
    return sum((obs.get(l, 0) - ref_dist[l]*n)**2 / (ref_dist[l]*n)
                for l in ALPHA if ref_dist[l] > 0)

ciphertext = k.op_vigenere_tr(target_plaintext, encrypt=True)
streams = [[] for _ in range(PERIOD)]
pos = 0
for c in ciphertext:
    if c in ALPHA:
        streams[pos % PERIOD].append(c)
        pos += 1

recovered = []
for stream in streams:
    best_shift, best_score = None, float('inf')
    for shift in range(10):  # shift restricted to a single digit
        decoded = [ALPHA[(ALPHA.index(c) - shift) % len(ALPHA)] for c in stream]
        score = chi_square(decoded, ref_dist)
        if score < best_score:
            best_score, best_shift = score, shift
    recovered.append(best_shift)

print("Recovered key:", "".join(map(str, recovered)))
\end{lstlisting}

\subsection{Known-Plaintext Brute Force (Section~\ref{sec:brute})}
\begin{lstlisting}[language=Python]
import kod8_engine as k
import time, random

def try_key(key_str, ciphertext, cipher_id):
    k.KOD8_KEY = key_str
    _SBOX = list(range(256))
    _seed = sum(int(c) * (i + 1) for i, c in enumerate(key_str))
    for _i in range(255, 0, -1):
        _seed = (_seed * 1103515245 + 12345) & 0x7FFFFFFF
        _j = _seed % (_i + 1)
        _SBOX[_i], _SBOX[_j] = _SBOX[_j], _SBOX[_i]
    k._SBOX = _SBOX
    k._SBOX_INV = [0]*256
    for _i, _v in enumerate(_SBOX):
        k._SBOX_INV[_v] = _i
    eng = k.Kod8(cipher_id=cipher_id)
    return eng.decrypt(ciphertext, cipher_id)

plaintext = "ATTACK AT DAWN TOMORROW"
enc, cid = k.Kod8(cipher_id="CL1").encrypt(plaintext)

random.seed(1)
N = 2000
trial_keys = [f"{random.randint(0, 99999999):08d}" for _ in range(N)]
trial_keys[N // 2] = "86247931"  # true key, inserted to confirm detection

t0 = time.perf_counter()
found = None
for i, cand in enumerate(trial_keys):
    if try_key(cand, enc, "CL1").strip("\x00") == plaintext.upper():
        found = (i, cand)
t1 = time.perf_counter()

per_trial = (t1 - t0) / N
print(f"{per_trial*1e6:.2f} us/trial, found at {found}")
print(f"Full 10^8 search estimate: {per_trial * 10**8 / 3600:.2f} hours")
\end{lstlisting}

\end{document}